\documentclass[12pt, a4paper]{article}

% --- Pacchetti Scientifici e Formattazione ---
\usepackage{amsmath, amssymb, amsfonts, amsthm}
\usepackage{booktabs, graphicx}
\usepackage[expansion=false]{microtype}
\usepackage[document]{ragged2e}
\usepackage{xcolor}

% --- Codifica e Font ---
\usepackage[utf8]{inputenc}
\usepackage[T1]{fontenc}
\usepackage[italian]{babel}
\usepackage{mathptmx}

% --- Gestione Margini e Layout ---
\usepackage[includeheadfoot]{geometry}
\geometry{
    a4paper,
    top=3cm,
    bottom=3cm,
    left=3.5cm,
    right=3.5cm,
}

% --- Gestione Interlinea e Figure ---
\usepackage{setspace}
\setstretch{1.3}
\usepackage[figuresright]{rotating}
\usepackage[pdftex, hidelinks]{hyperref}


\usepackage{subfiles}

% Document metadata
\title{Riassunto dei capitoli proposti della tesi di laurea.}
\author{Mattia Cavina}
\date{05 luglio 2026}

\begin{document}
\maketitle
\newpage

% --- Table of Contents ---
\tableofcontents

% --- Abstract---
\begin{abstract}
    la tesi verte sulla progettazione di un nuovo CPQ (Configure, Price, Quote)
    per l'azienda \textit{Cepi S.p.A.}, con l'obiettivo di sostituire il vecchio programma
    sviluppato in Microsoft Access.
\end{abstract}

% --- Situazione As-Is ---
\section{Situazione As-Is}
\subsection{Situazione Infrastrutturale Aziendale}
\subsection{Attuali software utilizzati}
\subsection{Organizzazione Aziendale e Flussi di lavoro}

% --- Interviste per la raccolta requisiti ---
\section{Raccolta dei requisiti}
\subsection{Scelta dei candidati}
\subsection{Modalità di svolgimento raccolta informazioni}
\section{Riassunto punti emersi per reparto}

% --- Analisi dei requisiti emersi dalle interviste ---
\section{Analisi dei requisiti}
\subsection{Requisiti funzionali}
\subsection{Requisiti non funzionali}

\subsection{Requisiti futuri shuld-have}

\section{Scelte infrastrutturali}
\subsection{Scelta dei deployment}
\subsection{Diagramma dei Deployment}
\subsection{Scelta dei Componenti}
\subsubsection{Diagramma dei Componenti}

% --- Progettazione del database ---
\section{Progettazione del database}
\subsection{Creazione di un primo modello concettuale}
\subsubsection{Esportazione Vecchio Db}
\subsubsection{Creazione regole di denominazione}

\subsection{Ristrutturazione database}
\subsubsection{Eliminazione ridondanze e normalizzazione}
\subsubsection{Creazione tabelle di lookup}
\subsubsection{Aggiunta tabelle non presenti}
\subsubsection{Creazione tabelle di Back Up}

\subsection{Schema logico}
\end{document}